
\section{代码示例}

本章展示如何在文章中插入代码块。以下是代码块的
使用示例，请根据实际需求修改代码内容。

\subsection{基础代码示例}

在此输入引出代码的过渡性文字，说明代码的功能和用途。

\begin{lstlisting}[language=C, caption={结构体定义示例}, label={lst:struct_example}]
#include <stdio.h>

// 结构体定义示例
typedef struct {
    char name[50];          // 名称字段
    int value;              // 数值字段
    float data;             // 数据字段
} ExampleStruct;

int main() {
    ExampleStruct example = {
        .name = "示例",
        .value = 100,
        .data = 3.14
    };
    
    printf("名称: %s\n", example.name);
    printf("数值: %d\n", example.value);
    
    return 0;
}
\end{lstlisting}

\subsection{函数定义示例}

在此输入函数相关的说明文字，介绍函数的设计思路
和使用方法。

\begin{lstlisting}[language=C, caption={函数示例代码}, label={lst:function_example}]
#include <stdio.h>

// 函数1：功能说明
void function_one() {
    printf(">>> 步骤1: 执行操作\n");
}

// 函数2：功能说明
void function_two() {
    printf(">>> 步骤2: 执行操作\n");
}

// 主函数
int main() {
    printf("=== 程序开始 ===\n\n");
    
    function_one();
    function_two();
    
    printf("\n=== 程序结束 ===\n");
    
    return 0;
}
\end{lstlisting}

\subsection{本章小结}

在此总结本章展示的代码示例及其使用要点。
